This section introduces the numerical formulation and physical models of LESGO and provides a brief overview of the OpenACC programming model.

\subsection{Pseudo-spectral LES formulation}
\label{sec:bg-numerics}

LESGO solves the filtered incompressible Navier--Stokes equations in
rotational form for high-Reynolds-number ABL flow,
%
\begin{equation}
\frac{\partial \tilde{u}_i}{\partial t}
+ \tilde{u}_j\!\left(\frac{\partial \tilde{u}_i}{\partial x_j}
- \frac{\partial \tilde{u}_j}{\partial x_i}\right)
= -\frac{\partial \tilde{p}^{*}}{\partial x_i}
- \frac{\partial \tau_{ij}}{\partial x_j}
+ f_i ,
\label{eq:les}
\end{equation}
%
where $\tilde{u}_i$ is the filtered velocity, $\tilde{p}^{*}$ is the 
modified pressure, $\tau_{ij}$ is the subgrid-scale (SGS) stress tensor, 
and $f_i$ represents body forces, including turbine forcing. The horizontal 
directions are periodic and discretized using Fourier collocation, whereas 
vertical derivatives are evaluated with second-order centered finite 
differences on a staggered grid. Dealiasing of nonlinear terms via the $3/2$ 
rule~\cite{orszag1971} requires forward and inverse FFTs across all velocity 
and vorticity components on a $3/2$-padded grid at every time step. 
Time integration employs a second-order Adams--Bashforth scheme, with
incompressibility enforced by a fractional-step projection. The pressure 
Poisson equation is diagonalized by 2D horizontal FFTs, yielding independent 
1D tridiagonal systems along $z$ for each horizontal wavenumber pair. At 
the lower boundary, an equilibrium log-law stress model~\cite{moeng1984} 
(or optionally an integral wall model~\cite{yang2015iwm}) closes the boundary 
condition. 

Consequently, FFT-based differentiation, dealiased products, and the spectral 
Poisson solve render the computational profile heavily FFT-dominated, 
in contrast to finite-difference atmospheric codes where FFTs are confined to 
the pressure solve~\cite{esclapez2025dales,vanheerwaarden2017}.
Prior GPU accelerations for such FFT-heavy solvers rely on CUDA Fortran,
including ports of finite-difference codes such as CaNS~\cite{costa2018,costa2021}  
and AFiD~\cite{zhu2018afid}, or custom asynchronous redesigns of 
pseudo-spectral codes~\cite{ravikumar2019}.

\subsection{Dynamic SGS and actuator-line turbine models}
\label{sec:bg-physics}

Two physics components dominate both the scientific fidelity and the
computational complexity of wind-farm LES in LESGO.

\textbf{Scale-dependent Lagrangian dynamic SGS (LASD).} An eddy
viscosity closes the SGS stress in Eq.~\ref{eq:les}. The model computes
its coefficient dynamically~\cite{germano1991}, averages it along
fluid-particle trajectories~\cite{meneveau1996}, and recovers scale
dependence through a second test
filter~\cite{porteagel2000,bouzeid2005}. Each
coefficient update applies two spectral test filters to nine tensor
fields (requiring batched 2-D FFTs across all vertical levels) and
advects four accumulator fields along particle back-trajectories via
trilinear interpolation. As the accuracy backbone of LESGO for
heterogeneous ABL flows, this model relies on an irregular memory access
pattern that intertwines filtering, history tracking, and interpolation.
Consequently, mapping this formulation to GPUs presents distinct
challenges compared to conventional local closures that rely
primarily on stencil operations~\cite{esclapez2025dales}, and it
has no counterpart in prior directive-based LES ports.

\textbf{Actuator-line turbine model (ATM).} Each turbine blade is
represented by discrete actuator points that sample the local LES
velocity via trilinear interpolation, look up tabulated airfoil
aerodynamic polars, and project the resulting forces back onto the grid
through a Gaussian convolution kernel~\cite{sorensen2002,martinez2015}.
Simulating multi-turbine farms introduces thousands of moving, off-grid
coupling points that interact with the flow field twice per time step
during velocity sampling and force projection (Section~\ref{sec:opt-atm}). 
Prior GPU-accelerated wind-farm simulation frameworks have largely 
focused on the ExaWind stack~\cite{sharma2024exawind,mullowney2021,kuhn2025amrwind}, 
which comprises second-order finite-volume codes tailored for 
geometry-resolved engineering simulations of finite wind farms. 
In contrast, LESGO addresses a distinct physical regime. Its 
pseudo-spectral discretization introduces zero numerical dissipation, 
enabling turbulence spectra, wake-recovery entrainment, and 
farm-boundary-layer interactions to be resolved without scheme-induced 
contamination. Furthermore, its horizontally periodic formulation 
directly models the fully developed wind-farm boundary layer that 
forms the foundation for atmospheric parameterizations. 

\subsection{GPU acceleration with OpenACC}
\label{sec:bg-openacc}

Early GPU acceleration efforts in atmospheric modeling relied either on 
converting select kernels to CUDA~\cite{michalakes2008,norman2015} or 
on complete solver rewrites~\cite{schalkwijk2012,vanheerwaarden2017}. 
In contrast, directive-based porting has emerged as the practical 
standard for legacy Fortran codebases~\cite{lapillonne2014,govett2017nim,giorgetta2022,esclapez2025dales,otero2019nek}.
OpenACC offloads execution by annotating loop nests with directives
(\texttt{parallel loop}), without altering the underlying source language.
Directives such as \texttt{create} and \texttt{copyin}/\texttt{copyout} 
handle explicit allocation and data movement, \texttt{update} synchronizes 
data on demand, and \texttt{declare create} maintains array residency 
throughout the program lifetime. Kernels compiled with \texttt{default(present)} 
enforce that operand arrays are already resident on the device.
Additionally, the \texttt{async} clause maps operations to asynchronous 
queues to overlap kernel execution, transfers, and library calls, 
with \texttt{wait} clauses ensuring correct data dependency ordering.

Although offloading individual compute loops via OpenACC follows 
well-established conventions, systematically managing data movement 
across a complex multiphysics solver remains the dominant performance 
bottleneck, which we address in Section~\ref{sec:method}.
